\documentclass[12pt,a4paper,oneside]{article}
\usepackage[brazil]{babel}
\usepackage[utf8]{inputenc}
\usepackage{olymp}

\setcounter{problem}{0}

\begin{document}

\begin{problem}{Substring}{substring}{2}
 
Uma \textit{substring} de uma \textit{string} $S$ é uma string $W$ tal que os caracteres de $W$ formem uma sequência existente em $S$. Uma substring tem uma posição inicial $P$ em $S$, e um comprimento $L$. Considere a seguinte string $S$:

\begin{center}
\begin{tabular}{|c|c|c|c|c|c|}
\hline
     0 & 1 & 2 & 3 & 4 & 5 \\
     B & A & T & M & A & N \\
\hline
\end{tabular}
\end{center}

A string ``BATMA'' é uma substring de $S$ com início na posição $P=0$ e comprimento $5$. A string ``MAN'' é uma substring com início na posição $P=3$ e comprimento $3$.

Faça um programa que, dada uma string $S$, uma posição $P$ e um comprimento $L$, imprima a substring de $S$ a partir da posição $P$ de comprimento $L$.

%\Notes
%
%Observações, se tiver.

\InputFile

A entrada é composta por uma string $S$ e dois inteiros $P$ e $L$, representando a posição inicial da substring e seu comprimento, respectivamente.

Todas as palavras são compostas apenas por letras maiúsculas, sem caraceteres especiais. Restrições: $1 \leq len(S) \leq 100$, $0 \leq P < len(S)$, $1 \leq L \leq len(S) - P$, onde $len(S)$ é o comprimento da string $S$.

\OutputFile

Seu programa deve imprimir a substring solicitada na entrada, em letras maiúsculas.

% \Notes
% Preste atenção aos tipos das variáveis, pois $F(90) \approx 2^{61}$.

\Example

\begin{example}
\exmp{
MANGA 0 3
}{
MAN
}%
\end{example}

\begin{example}
\exmp{
PARALELEPIPEDO 4 4
}{
LELE
}%
\end{example}

\begin{example}
\exmp{
VINGADORES 3 4
}{
GADO
}%
\end{example}

\begin{example}
\exmp{
PROGRAMACAO 0 11
}{
PROGRAMACAO
}%
\end{example}

%Comentários sobre o exemplo, se necessário.

\end{problem}

\end{document}
